% PraTeX の plain format で Vaak bridge を使う最小文書。
% 実行例:
%   pratex "&plain" examples/plain-vaak.tex

% \vaakdef は定義時に Vaak program を検査・組み立てし、呼出し時には再利用する。
% let と var は \directvaak / \vaakdef の中で通常どおり使える。
\count5=21
\count6=1
\vaakdef\vaakanswer{
  let base : i32 := count[5];
  var doubled : i32 := base * 2;
  doubled + count[6]
}
\count100=\vaakanswer

% count と dimen は起動時から host 値として見えている。
% 直接 count[5] と書けるが、&= で明示的な host alias を作ることもできる。
\count101=\directvaak{
  var registers : i32 array alias &= count;
  let previous : i32 := registers[5];
  registers[5] := previous + 1;
  registers[5]
}

\message{[plain-vaak answer=\the\count100;alias=\the\count101;host=\the\count5]}

\shipout\vbox{
  \hbox{PraTeX plain format with Vaak}
  \vskip 6pt
  \hbox{vaakdef result: \the\count100}
  \hbox{host alias result: \the\count101}
  \hbox{host count[5] after writeback: \the\count5}
}
\end
